\section{What a process-based model is made of}
A model of the soil--plant--atmosphere system, such as Daisy, is built from a few kinds of element, and the process map of Figure~\ref{fig:dmap} shows all of them. The \kw{state variables} are the stores that the model keeps from one time step to the next (the water content of each layer, the temperature, the nitrogen of each pool, the dry matter of each organ). The \kw{processes} are the flows that change the state (infiltration, uptake, mineralisation). The \kw{parameters} are the constants of the processes (the conductivity, the rate of decay), and the \kw{forcing} is the input from outside that the model does not compute (the weather, the management). At every step the model reads the forcing, computes the flows from the state and the parameters, and updates the state.

\section{The principle of conservation}
All the boxes of the map are linked by one rule: \kw{what enters a store and does not leave it, stays in it}. For any store $S$,
\begin{equation*}\boxed{\frac{\mathrm dS}{\mathrm dt}=\sum F_{in}-\sum F_{out}}\end{equation*}
the change of storage is the sum of the inflows less the sum of the outflows. This is the law of conservation of mass (or of energy, for the heat of the soil), and it is the reason why a model of this kind must close its balance. In the map, a solid arrow is a flow of matter, so that a quantity that leaves one box must arrive at another or leave the system; a dashed arrow is an exchange of information (a control), which does not carry mass.

\section{Time step and scale}
The processes of the map have different speeds. The infiltration of a storm takes minutes, the redistribution hours, the growth of the crop days to months and the decay of the soil organic matter years. The model uses a \kw{time step} short enough for the fastest process that matters (an hour for the water flow, a day for the crop) and runs the slow processes with the same step, which costs little. The \kw{spatial} scale is the soil column (a point in the field) divided into layers, thin near the surface where the gradients are large and thicker at depth; the results are for a field of uniform soil and management.

\section{The interactions among the parts of the model}
The parts of the map affect each other in both directions, which is what makes the system non-linear. The water of the soil affects the crop (uptake, stress) and the crop affects the water (transpiration, interception); the temperature of the soil affects the biology (decay, nitrification) and the biology produces heat; the nitrogen of the soil limits the growth, and the growth takes the nitrogen. These \kw{feedback loops} are the subject of the title of this book. A change made at one place (a mulch that reduces evaporation) is followed through the whole map and, as Chapter~\ref{chap:09} shows, the result is not always the one expected.

\section{Reading the map: a method}
To follow a question on the map, (1) find the box where the quantity lives, (2) follow the solid arrows into it to find its inputs and out of it for its outputs, (3) follow the dashed arrows into the process boxes to find what controls the rates, and (4) find the log column that holds the quantity for a run, so that the answer can be checked against the output. The sections of this chapter and of the two chapters that follow use this method, box by box.
